\GXNSFBodyText

\noindent\textbf{研究内容}\par

\noindent\textbf{研究内容一：跨媒介职业事件链构建}\par
围绕佐佐木希公开职业资料建立演示事件链，记录时间、平台、内容类型、角色身份、事件属性和可见反馈。数据只纳入官方网站、作品页面、公开新闻稿和可依法使用的平台元数据，不采集私人账号、非公开互动或敏感个人信息。针对日文、中文和英文资料，建立实体对齐、时间归一化、重复事件消解与来源置信度规则，形成可追溯的多语言事件表示。

\noindent\textbf{研究内容二：多模态职业状态识别}\par
将文本主题、视觉风格、内容更新节奏和平台网络位置映射到统一状态空间。先由平台专属校准映射 $g_p$ 将原始观测 $\bm{x}_t^{\mathrm{raw}}$ 转换为共享向量 $\bm{x}_t=g_p(\bm{x}_t^{\mathrm{raw}})$。设隐含职业状态为 $z_t$，平台环境与外部扰动为 $\bm{e}_t$，状态数为 $K$，则从状态 $j$ 向候选状态 $k$ 的条件转移概率可表示为
\begin{equation}
  p(z_t=k\mid z_{t-1}=j,\bm{x}_t,\bm{e}_t)=
  \frac{\exp\{a_{jk}+\bm{b}_k^{\mathsf T}\bm{x}_t+\bm{c}_k^{\mathsf T}\bm{e}_t\}}
  {\sum_{h=1}^{K}\exp\{a_{jh}+\bm{b}_h^{\mathsf T}\bm{x}_t+\bm{c}_h^{\mathsf T}\bm{e}_t\}}.
  \label{eq:gx-state-transition}
\end{equation}
其中 $a_{jk}$ 表示基础转移倾向，$\bm{b}_k$ 与 $\bm{c}_k$ 分别表示校准后观测和环境系数；分母对全部候选当期状态归一化。本项目采用条件状态模型，由 $g_p$ 统一校准平台差异。通过条件马尔可夫基线、变点检测与时序表示学习的对照实验，识别若干候选状态。阶段名称依据模型输出及其观测证据加以解释，而不预先设定。

\noindent\textbf{研究内容三：扰动响应与职业韧性估计}\par
以公开可核验的重大职业节点、平台切换或外部舆论冲击为候选扰动，比较其前后内容供给、媒介多样性、互动结构和语义一致性的变化。职业节点与平台切换首先进入描述性时序比较；只有时点唯一性、对照可比性和前趋势等条件得到支持的事件，才进入准实验稳健性分析。先用扰动前固定窗口估计基线 $B$，再在长度固定的扰动后窗口 $T$ 中计算四个归一化分量，定义综合韧性指标
\begin{equation}
  R_T=w_S S_T+w_D D_T+w_C C_T+w_Q Q_T,
  \label{eq:gx-resilience}
\end{equation}
其中 $S_T$ 表示相对基线的恢复速度，$D_T$ 表示媒介多样性，$C_T$ 表示核心叙事一致性，$Q_T$ 表示波动稳定性。权重均非负且总和为 1。四个分量均映射到 $[0,1]$，数值越大表示该维度韧性越强；未在窗口 $T$ 内恢复到预设基线容差时，$S_T$ 记为 0。主分析采用等权并同时报告四个分量，替代权重只用于敏感性分析；恢复窗口之后的独立效标仅用于检验指标效度，不参与 $R_T$ 计算。

\noindent\textbf{研究内容四：跨案例验证与组件适用性测试}\par
选取公开资料较完整、职业跨度相近的东亚女性艺人作为匿名化对照组，检验完整状态模型是否过度贴合佐佐木希单一个案。同时，以“单一文化内容 IP 或单一区域品牌—固定观察窗”为分析单位构造广西—东盟小规模样本，考察多语言实体对齐、平台校准和扰动窗口等方法组件在区域文化传播材料中的适用性。两类样本分别服务于完整模型验证与组件测试，比较结果用于识别方法边界，而非评价人物或品牌优劣。

\noindent\textbf{研究目标}\par

\noindent\textbf{总体目标}\par
构建一个具有来源可追溯、状态可解释、结果可复核特征的跨媒介职业韧性分析框架，刻画多层媒介网络中“平台迁移—身份重组—扰动恢复”的时序关联。

\noindent\textbf{具体目标}\par
\begin{enumerate}[label=（\arabic*）,leftmargin=*,itemsep=0pt]
  \item 建立多语言公开职业事件的统一编码规范与质量控制流程；
  \item 提出能够区分平台差异与主体状态变化的多模态时序模型；
  \item 给出职业韧性的操作化定义、估计方法及适用边界；
  \item 通过同构人物对照检验完整模型稳健性，并在广西—东盟样本上测试方法组件的探索性适用性。
\end{enumerate}

\noindent\textbf{拟解决的关键科学问题}\par

\noindent\textbf{多平台异质观测如何映射到可比较的隐状态}\par
不同媒介的更新频率、内容长度和反馈指标并不处于同一量纲。在保留平台特性的前提下，哪些观测信息能够分离平台测量偏差与主体状态变化，并形成跨时间、跨平台可比较的状态表示？

\noindent\textbf{外部扰动前后跨媒介状态路径呈现何种差异}\par
同一扰动前后可能出现短期热度升高、内容停滞或平台迁移。扰动强度、平台组合与既有职业资本在什么条件下对应不同的时序路径？系统回到原有状态与进入新稳定状态，分别呈现哪些可观测特征？

\noindent\textbf{单案例发现如何获得可反驳的同构检验}\par
如何建立从个案发现到同构人物对照的证据链，区分佐佐木希案例的特定历史条件与可迁移的状态表征？当模型在对照样本中失效时，哪些条件构成其适用边界？区域数字文化场景则用于检验通用方法组件，不延伸为人物职业机制的验证。
